\clearpage\thispagestyle{empty}
\pdfbookmark[0]{Source for the Report on Vesta, Juno, Pallas and Ceres}{planets1808july-sources}
\begin{center}\textbf{Source for the Report on Vesta, Juno, Pallas and Ceres}\end{center}
\noindent This AI translation follows the report dated 4 July 1808 in \href{https://gdz.sub.uni-goettingen.de/id/PPN236007467}{Gauss’s \textit{Werke}, vol.~VI} (Göttingen, 1874), pp.~303--304. It appeared under the heading \textit{Göttingen} in the combined issues 107 and 108 of the \textit{Göttingische gelehrte Anzeigen}, \href{https://manifests.sub.uni-goettingen.de/iiif/image/gdz:PPN319721507_1808:00001071/full/1600,/0/default.jpg}{pp.~1065}, \href{https://manifests.sub.uni-goettingen.de/iiif/image/gdz:PPN319721507_1808:00001072/full/1600,/0/default.jpg}{1066}, \href{https://manifests.sub.uni-goettingen.de/iiif/image/gdz:PPN319721507_1808:00001073/full/1600,/0/default.jpg}{1067} and \href{https://manifests.sub.uni-goettingen.de/iiif/image/gdz:PPN319721507_1808:00001074/full/1600,/0/default.jpg}{1068}. The collected edition cites issue 107. The report refers to Gauss in the third person.
\par\medskip
The first printing gives two observations of Juno on p.~1067. The collected edition directs readers to those rows in the letter of 27 June. The report’s Vesta table retains its perihelion entry, its own precision for the node and inclination, and the irregular separators of the collected edition. Translation and typesetting produced by AI.
\clearpage
